\section{Distribution definitions and dataset construction}
\label{app:distributions}
All target laws use the following parameter conventions. Densities are zero
outside the stated support; discrete formulas are probability masses.
\begin{align}
 \text{Normal: }\quad
 p(x)&=\frac{1}{\sigma\sqrt{2\pi}}
       \exp\!\left[-\frac{(x-\mu)^2}{2\sigma^2}\right],
       &&x\in\R,\\
 \text{Uniform: }\quad p(x)&=\frac{1}{b-a},&&a\leq x\leq b,\\
 \text{Exponential: }\quad p(x)&=\lambda e^{-\lambda x},&&x\geq0,\\
 \text{Beta: }\quad p(x)&=\frac{x^{a-1}(1-x)^{b-1}}{\mathrm B(a,b)},
       &&0<x<1,\\
 \text{Gamma: }\quad p(x)&=\frac{x^{k-1}e^{-x/s}}{\Gamma(k)s^k},&&x>0,\\
 \text{Binomial: }\quad p(x)&=\binom{n}{x}\rho^x(1-\rho)^{n-x},
       &&x\in\{0,\ldots,n\},\\
 \text{Weibull: }\quad p(x)&=\frac{k}{s}\left(\frac{x}{s}\right)^{k-1}
                e^{-(x/s)^k},&&x>0,\\
 \text{Poisson: }\quad p(x)&=e^{-\lambda}\frac{\lambda^x}{x!},
       &&x\in\{0,1,\ldots\}.
\end{align}
Here $\mathrm B(a,b)=\Gamma(a)\Gamma(b)/\Gamma(a+b)$. Gamma and Weibull
use shape/scale conventions, while exponential and Poisson use rates.
Support endpoints with zero continuous probability do not affect sampling.

Parameters are drawn independently and uniformly from the ranges in
Table~\ref{tab:ranges}, then rounded to five decimals. For binomial $n$,
the integer ranges are inclusive. Training and validation use the same
range rules but disjoint exact parameter combinations. The held-out
family ranges are used only for their test prompts.

\begin{table}[htbp]
\centering\small
\caption{Parameter-generation ranges. $U[a,b]$ denotes a continuous uniform
draw before parameter rounding. The uniform family's width is $w=b-a$.}
\label{tab:ranges}
\begin{tabular}{@{}lp{0.37\linewidth}p{0.37\linewidth}@{}}
\toprule
Family & Training / validation / new wording & Unseen parameters or families\\
\midrule
Normal & $\mu\sim U[-5,5]$, $\sigma\sim U[0.3,3]$ &
 $\mu\sim U[6,9]$, $\sigma\sim U[3.5,5]$\\
Uniform & $a\sim U[-5,2]$, $w\sim U[0.5,8]$ &
 $a\sim U[4,7]$, $w\sim U[9,12]$\\
Exponential & $\lambda\sim U[0.2,3]$ & $\lambda\sim U[3.5,5]$\\
Beta & $a,b\sim U[0.5,5]$ & $a,b\sim U[6,9]$\\
Gamma & $k\sim U[0.5,5]$, $s\sim U[0.3,3]$ &
 $k\sim U[6,9]$, $s\sim U[3.5,5]$\\
Binomial & $n\in\{5,\ldots,50\}$, $\rho\sim U[0.1,0.9]$ &
 $n\in\{55,\ldots,80\}$, $\rho\sim U[0.1,0.9]$\\
Weibull & Not used & $k\sim U[0.6,3]$, $s\sim U[0.5,4]$\\
Poisson & Not used & $\lambda\sim U[0.5,20]$\\
\bottomrule
\end{tabular}
\end{table}

The dataset seed is 20261117. The three prompt templates are:
\par\noindent\begin{minipage}{\linewidth}
\begin{quote}\small
\textbf{Training template 0:} Sample exactly one number from [distribution
description]. Output only the number.\newline Answer:\par
\textbf{Training template 1:} Generate one random draw from [distribution
description]. Return only its numeric value.\newline Answer:\par
\textbf{Held-out template 2:} I need a single independent observation drawn
according to [distribution description].\newline Number:
\end{quote}
\end{minipage}\par
Each label is followed by a newline. Descriptions spell out the parameter
names, for example ``a normal distribution with mean 7.83271 and standard
deviation 3.98876.'' Exact prompts, rather than this illustrative parameter
example, are stored in the committed test-case file. Training-family
selection cycles through six families; test-category selection cycles
through the three conditions while avoiding a coupling that would omit
families. Validation contains 160 configurations (1,280 target answers for
SFT); the RL probes use its first 12 configurations, two from each family.

\section{Gradient and objective details}
\subsection{Quantile bins and invalid responses}
\label{app:bins}
With a numeric-bin budget $K_{\max}=32$, begin with internal quantiles
$F_c^{-1}(j/32)$, $j=1,\ldots,31$. For a continuous law, round each to
the output grid and place the bin boundary a half-grid above it:
\begin{equation}
 e_j=\operatorname{round}_{\delta}\!\left(F_c^{-1}(j/32)\right)+\delta/2.
\end{equation}
Remove duplicate boundaries, CDF values at 0 or 1, and boundaries giving
zero mass. Numeric-bin masses are successive differences of
$[0,F_c(e_1),\ldots,F_c(e_{K-1}),1]$. These boundaries separate rounded
grid outputs, so CDF differences account for the rounded target mass.

For a discrete law with a finite support of at most 32 integers, all
adjacent support values are separated by half-integer boundaries. For
larger or infinite discrete supports, unique quantiles plus $1/2$ define
boundaries, followed by the same CDF-difference construction. Bin
assignment uses right-sided boundary search. A finite parsed number is
invalid for the reward if it is outside support (with half-grid tolerance
for continuous laws), or noninteger for a discrete law. Malformed text
and unfinished completions are also invalid.

The penalty for invalid mass can be seen at the population histogram
level. If $q_\bot<1$ and $\overline q_j=q_j/(1-q_\bot)$ normalizes the
numeric bins, then
\begin{equation}
 \KL(P\Vert Q)=\KL(P\Vert\overline Q)-\log(1-q_\bot).
 \label{eq:invalidpenalty}
\end{equation}
Thus ignoring the zero-target-mass term does not mean ignoring invalid
outputs; the loss still penalizes the probability diverted away from
valid bins. In implementation, the same normalization effect operates on
the smoothed empirical histogram.

\subsection{Unbiasedness of the leave-one-out policy gradient}
Fix a prompt $c$ and let the joint sampling law be
$\prod_i\pi_\theta(Y_i\mid c)$. Differentiating the expected batch reward
gives
\begin{align}
 \nabla_\theta\E R_N
 &=\sum_{Y_{1:N}} R_N\,
   \nabla_\theta\prod_i\pi_\theta(Y_i\mid c)\\
 &=\E\!\left[R_N\sum_i\nabla_\theta\log\pi_\theta(Y_i\mid c)\right].
\end{align}
For $b_i=R_{N-1}(c,Y_{-i})$, independence conditional on the prompt implies
\begin{align}
 \E[b_i\nabla_\theta\log\pi_\theta(Y_i\mid c)]
 &=\E_{Y_{-i}}\!\left[b_i\,
       \E_{Y_i}\nabla_\theta\log\pi_\theta(Y_i\mid c)\right]\\
 &=\E_{Y_{-i}}[b_i\nabla_\theta 1]=0.
\end{align}
Although the baseline's distribution changes with $\theta$, it is used as
a stop-gradient control variate; no baseline derivative is included.
This establishes Equation~\eqref{eq:groupgradient}. It relies on fresh
on-policy samples, consistent actor likelihoods, and the stated
conditional independence assumption. Reusing the same samples for multiple
unadjusted actor steps would break this on-policy justification.

\subsection{Closed-form leave-one-out advantage}
Write $j_i$ for the bin of action $i$, $n_{j_i}\geq1$ for its group count,
and $B$ for all categories, including invalid. Expanding the reward as
\begin{equation}
 R_N=-\sum_{j:p_j>0}p_j\log p_j
       +\sum_{j:p_j>0}p_j\log(n_j+\alpha)-\log(N+\alpha B)
\end{equation}
shows that removing one action only changes its selected count and the
normalizer. Therefore
\begin{equation}
 R_N-R_{N-1}=
 \log\frac{N-1+\alpha B}{N+\alpha B}
 +p_{j_i}\log\frac{n_{j_i}+\alpha}{n_{j_i}-1+\alpha}.
 \label{eq:closedadvantage}
\end{equation}
The implementation computes all advantages by this expression rather
than recomputing $N$ full histograms. If the action is invalid,
$p_{j_i}=0$ and its advantage is negative. Multiplying by $N$ and
averaging sequence losses recovers the group sum in
Equation~\eqref{eq:groupgradient}. Unit tests compare this estimator
against exact outcome enumeration on a small categorical group.

\subsection{What finite-sample histogram optimization does and does not imply}
\label{app:finite}
Let $q_j$ be the true model probability of histogram category $j$ at a
fixed prompt. Each marginal count $M_j$ has a binomial distribution
$\operatorname{Binomial}(N,q_j)$, even though the complete count vector is
multinomial. The expected reward is exactly
\begin{equation}
 \E R_N=-\sum_{j:p_j>0}p_j\log p_j-\log(N+\alpha B)
       +\sum_{j:p_j>0}p_j\E_{M_j}\log(M_j+\alpha).
 \label{eq:expectedreward}
\end{equation}
The expectation of a log count differs from the log of its expectation.
Consequently, maximizing the finite-$N$ reward is not identical to
minimizing $\KL(P\Vert Q)$, and its optimum need not be exactly $Q=P$.
The estimator's unbiased policy gradient is for this finite-sample
objective. As $N$ grows with fixed $B$ and $\alpha$, the empirical
smoothed histogram converges to $Q$; if all positive-target bins have
positive model probability, the realized KL converges to population bin
KL. This remains a statement about bins, not within-bin densities.

For a continuous target, an unsmoothed empirical measure places mass on a
finite set whereas the target has density over an interval. Directly
treating that empirical measure as a continuous density makes the
forward KL ill-defined or infinite. Explicit binning supplies the common
finite measurable partition needed by our implemented reward. A kernel
density alternative would introduce bandwidth and tail assumptions; we
do not use one here.

\subsection{Reference-KL score-function derivative}
\label{app:reference}
For fixed reference policy $\pi_0$ on the same capped completion space,
\begin{align}
 K(\theta)&=\sum_y\pi_\theta(y\mid c)
                  \log\frac{\pi_\theta(y\mid c)}{\pi_0(y\mid c)},\\
 \nabla_\theta K
 &=\E_{y\sim\pi_\theta}
   \left[\left(\log\frac{\pi_\theta(y\mid c)}{\pi_0(y\mid c)}+1\right)
            \nabla_\theta\log\pi_\theta(y\mid c)\right]\\
 &=\E_{y\sim\pi_\theta}
   \left[\log\frac{\pi_\theta(y\mid c)}{\pi_0(y\mid c)}
            \nabla_\theta\log\pi_\theta(y\mid c)\right].
\end{align}
The last step uses zero expected policy score. This is why
Equation~\eqref{eq:surrogate} subtracts the detached log-ratio from the
reward coefficient. Simply backpropagating an on-policy sampled log
ratio while treating sampled paths as fixed would omit the distribution's
score contribution. Both adapters share frozen base weights, while the
reference LoRA weights stay frozen. Disabling dropout for generation and
likelihood recomputation is necessary to score the sampled policy
consistently in this implementation.

\subsection{Direct soft-target KL: implemented, not a completed experiment}
For completeness, the code also supports supervision from a finite
canonical target law $P^{\mathrm{fin}}$ with at most 1,001 output strings.
For continuous laws, a rounded grid spans underlying quantiles 0.001 to
0.999; CDF differences define masses and edge bins absorb outer tails.
Discrete support is represented or coarsened under the same point budget.
Each output's token sequence, including EOS, is inserted into a prefix
trie. If $s$ is a target prefix, define
\begin{equation}
 P^{\mathrm{fin}}(v\mid s,c)=
 \frac{\sum_{y:sv\preceq y}P^{\mathrm{fin}}(y\mid c)}
      {\sum_{y:s\preceq y}P^{\mathrm{fin}}(y\mid c)}.
\end{equation}
Along target-sampled paths, the expected sum of next-token forward KLs
satisfies the chain rule
\begin{equation}
 \E_{Y\sim P^{\mathrm{fin}}}\sum_{t=1}^{T(Y)}
 \KL\!\left(P^{\mathrm{fin}}(\cdot\mid Y_{<t},c)
          \Vert\pi_\theta(\cdot\mid Y_{<t},c)\right)
 =\KL(P^{\mathrm{fin}}\Vert\pi_\theta).
\end{equation}
The model's softmax normalizer includes the full vocabulary; it is not
renormalized to valid trie children. Invalid output mass is therefore
penalized. The code sums losses over prefixes and averages over sampled
paths. The default uses three paths per training configuration, giving
6,000 paths, with a 64-trie cache. This is teacher-supervised distribution
matching, not a reward from model-generated samples. Its finite target
also differs from the full rounded law used by SFT and evaluation. No
full-model soft-KL results are included in this paper.

\section{Extended numerical results}
\label{app:extended}
\begin{table}[htbp]
\centering\small
\caption{Paired prompt-bootstrap differences, RL minus SFT. Negative
differences favor RL. Intervals use 10,000 resamples of the 100 prompt
indices and percentile endpoints. They do not incorporate training-seed
variability or multiple-comparison correction.}
\begin{tabular}{@{}llrr@{}}
\toprule
Metric & Aggregate & Difference & Exploratory 95\% interval\\
\midrule
Histogram KL & Mean & -0.005107 & [-0.008401, -0.001804] \\
Histogram KL & Median & -0.013447 & [-0.021973, +0.000740] \\
KS distance & Mean & -0.005296 & [-0.008375, -0.002172] \\
KS distance & Median & -0.012387 & [-0.022164, -0.000011] \\
Normalized Wasserstein & Mean & -0.003660 & [-0.006756, -0.000684] \\
Normalized Wasserstein & Median & -0.010959 & [-0.016946, +0.002674] \\
\bottomrule

\end{tabular}
\end{table}

\begin{table}[htbp]
\centering\small\setlength{\tabcolsep}{5pt}
\caption{Response quality over 100,000 requests per checkpoint.
Violations are valid numeric answers outside support or noninteger for
a discrete target. Truncation is measured over all requests.}
\label{tab:quality}
\begin{tabular}{@{}lrrrrr@{}}
\toprule
Policy & Valid count & Validity & Violations & Violations / valid & Truncated\\
\midrule
Pretrained & 6,838 & 6.838\% & 1,359 & 19.874\% & 88.625\% \\
SFT & 98,677 & 98.677\% & 1,330 & 1.348\% & 0.648\% \\
SFT + RL & 98,591 & 98.591\% & 1,208 & 1.225\% & 0.674\% \\
\bottomrule

\end{tabular}
\end{table}

\begin{table}[htbp]
\centering\small
\caption{All recorded RL validation probes: 12 fixed configurations,
256 draws each. The final tested adapter is update 250. The best mean KL
occurs at update 225, whose adapter was retained but not test-evaluated.}
\label{tab:validation}
\begin{tabular}{@{}rrrrr@{}}
\toprule
Update & Mean histogram KL & Median KS & Median normalized $W_1$ & Validity\\
\midrule
0 & 0.070539 & 0.068078 & 0.039160 & 100.00\% \\
25 & 0.066361 & 0.070032 & 0.034210 & 100.00\% \\
50 & 0.072817 & 0.069250 & 0.039073 & 100.00\% \\
75 & 0.072290 & 0.061455 & 0.034566 & 100.00\% \\
100 & 0.069355 & 0.064310 & 0.028803 & 100.00\% \\
125 & 0.067104 & 0.060269 & 0.028083 & 100.00\% \\
150 & 0.066409 & 0.063268 & 0.033190 & 100.00\% \\
175 & 0.065829 & 0.063264 & 0.031405 & 100.00\% \\
200 & 0.067312 & 0.065217 & 0.029034 & 100.00\% \\
225 & 0.064350 & 0.070032 & 0.031264 & 100.00\% \\
250 & 0.065246 & 0.056670 & 0.030995 & 100.00\% \\
\bottomrule

\end{tabular}
\end{table}

All 100 prompt-level base/SFT/RL outcomes are included in the source
package's supplementary CSV. The paper uses the already-recorded test
metrics and the post-hoc histogram KL analysis; it does not select a
different checkpoint or decoding configuration to improve the results.
The repository additionally records canonical formatting, repeated-value
rates, sample means and standard deviations, support violations, and
metric overflow flags per prompt. These diagnostics are descriptive,
not a comprehensive test of random-number quality.

\FloatBarrier
\section{Artifacts and reproduction}
\label{app:records}
The numeric code and records are available at
\href{https://github.com/YashSolanki2007/llm-random-number-generator/tree/27fad3ec875afb2941f34ab873e8394ae24f25e8}{repository commit \texttt{27fad3e}}.
It includes all 200,000 SFT/RL test responses, per-prompt metrics,
training history, configs, and provenance hashes. The earlier SFT records
include base metrics and two cases' raw responses. Training rollouts and
adapters remain in W\&B artifacts and cluster storage. This manuscript's
source archive includes generated tables, vector figures, the bibliography,
and a supplementary per-prompt CSV.

The production environment records PyTorch 2.8.0+cu126, Transformers
4.57.1, PEFT 0.17.1, Accelerate 1.10.1, NumPy 2.2.6, SciPy 1.15.3,
and W\&B 0.31.0, with CUDA 12.6 on an NVIDIA H200. Production training
and evaluation require online W\&B logging. No credentials are present
in this manuscript or its source archive.

\paragraph{Run records.}
\begin{itemize}
\item SFT training:
\href{https://forge.coreweave.com/wandb/yash2007-university-of-washington/stochastic-llm/runs/xcj01bg1}{\texttt{xcj01bg1}}.
\item Base and SFT evaluation:
\href{https://forge.coreweave.com/wandb/yash2007-university-of-washington/stochastic-llm/runs/rfddft7y}{\texttt{rfddft7y}} and
\href{https://forge.coreweave.com/wandb/yash2007-university-of-washington/stochastic-llm/runs/ao4xvmp5}{\texttt{ao4xvmp5}}.
\item RL training and evaluation:
\href{https://forge.coreweave.com/wandb/yash2007-university-of-washington/stochastic-llm/runs/tx91cho0}{\texttt{tx91cho0}} and
\href{https://forge.coreweave.com/wandb/yash2007-university-of-washington/stochastic-llm/runs/n9d9aegn}{\texttt{n9d9aegn}}.
\item Final SFT/RL comparison:
\href{https://forge.coreweave.com/wandb/yash2007-university-of-washington/stochastic-llm/runs/lyb3xv9l}{\texttt{lyb3xv9l}}.
\end{itemize}
The RL timing benchmark, training, and evaluation were Slurm jobs
493753, 493754, and 493755, respectively; all completed successfully.
These identifiers are provenance, not instructions to reuse outputs.

\paragraph{Regenerate analysis from the repository.}
The following commands require no GPU and create no online training run:
\begin{quote}\small\ttfamily
python -m pip install -e '.[dev,analysis]'\\
python -m pytest -q\\
python scripts/analyze\_rl\_results.py --recorded\\
python paper/generate\_assets.py
\end{quote}
The manuscript compiles independently of the experiment data. Regenerating
its tables and figures requires the repository records.
